\documentclass[11pt,a4paper,oneside]{book}
\usepackage{mathscommun}
\hypersetup{pdftitle={Mathématiques de troisième - Cours et démonstrations}}
\begin{document}
\titrecollection{Cours et démonstrations}{Comprendre, justifier, puis savoir utiliser.}{Un cours organisé en vingt chapitres : définitions précises, preuves expliquées, exemples résolus et figures vectorielles. Les prérequis sont repris et les approfondissements sont signalés.}
\frontmatter
\chapter*{Comment travailler avec ce cours}
\addcontentsline{toc}{chapter}{Comment travailler avec ce cours}
Le but n'est pas de réciter une collection de formules. On cherche à comprendre ce que désignent les objets, pourquoi une propriété fonctionne, dans quelles conditions elle s'applique et comment la reconnaître dans un problème.

Pour chaque chapitre : lire les objectifs ; refaire les exemples sans regarder ; expliquer une démonstration à voix haute ; résoudre les quatre exercices correspondants du cahier ; consulter ensuite le corrigé. Espacer les reprises est plus utile que relire plusieurs fois immédiatement.

\begin{attention}[Ce que signifie « démontrer »]
Une définition fixe le sens d'un mot : elle ne se démontre pas. Un théorème se démontre à partir de définitions et de résultats déjà établis. Quelques principes de base sont admis et recensés dans le guide. Une figure ou une simulation aide à comprendre ; elle ne remplace pas une preuve générale.
\end{attention}

Le périmètre est le programme français de troisième en vigueur en 2026--2027, complété par les acquis du cycle 4 dont il dépend. Les notations d'ensembles, certaines identités, les inéquations et quelques preuves sont des compléments utiles, explicitement signalés. Ce volume ne prétend pas couvrir à lui seul le nouveau programme de troisième applicable à partir de septembre 2028.

La table de couverture détaillée, les sources officielles, la progression proposée, le registre des démonstrations et les consignes de compilation sont dans \emph{Guide, couverture et révisions}. Les cahiers d'exercices et de corrigés utilisent les mêmes numéros de chapitre.

\tableofcontents
\mainmatter
\subfile{chapitres/01_raisonner.tex}
\subfile{chapitres/02_nombres_fractions.tex}
\subfile{chapitres/03_arithmetique.tex}
\subfile{chapitres/04_puissances_racines.tex}
\subfile{chapitres/05_calcul_litteral.tex}
\subfile{chapitres/06_equations.tex}
\subfile{chapitres/07_proportionnalite.tex}
\subfile{chapitres/08_fonctions.tex}
\subfile{chapitres/09_fonctions_affines.tex}
\subfile{chapitres/10_statistiques.tex}
\subfile{chapitres/11_probabilites.tex}
\subfile{chapitres/12_geometrie_plane.tex}
\subfile{chapitres/13_pythagore.tex}
\subfile{chapitres/14_thales_similitude.tex}
\subfile{chapitres/15_trigonometrie.tex}
\subfile{chapitres/16_transformations.tex}
\subfile{chapitres/17_grandeurs_mesures.tex}
\subfile{chapitres/18_espace.tex}
\subfile{chapitres/19_algorithmique.tex}
\subfile{chapitres/20_outils_problemes.tex}
\end{document}
